\documentclass[10pt,a4paper]{amsart}
\usepackage[margin=19mm]{geometry}
\usepackage{amsmath,amssymb,mathtools}
\usepackage[scr=rsfs,cal=euler]{mathalfa}
\usepackage{xfrac}
\usepackage[unicode,hidelinks,bookmarks=false]{hyperref}
\newcommand{\ie}{i.e.\ }
\newcommand{\cf}{cf.\ }
\newcommand{\resp}{respectively\ }
\newcommand{\Subsection}{\subsection}
\providecommand{\nobreakdash}{\nobreak\hbox{-}}
\setlength{\emergencystretch}{2em}
\multlinegap0pt
\newcommand{\Mup}{M_{up}}
\usepackage{float}
\usepackage{mathtools}
\usepackage{tikz}
\newtheorem{theorem}{Theorem}
\title{A Combinatorial Framework for the Pons--Batle Identity: Young Tableaux, Lattice Paths, and Limit Laws}
\author{Hexuan Liu, Michael Wallner, Guan-Ru Yu}
\date{Source: arXiv:2605.07587v1 (CC BY 4.0)}
\begin{document}
\maketitle

\noindent\emph{Source and licence.} A Combinatorial Framework for the Pons--Batle Identity: Young Tableaux, Lattice Paths, and Limit Laws. Version arXiv:2605.07587v1 (CC BY 4.0), distributed by arXiv under the Creative Commons Attribution 4.0 International licence, http://creativecommons.org/licenses/by/4.0/, as recorded in the arXiv licence metadata for this exact version and shown on the abstract page https://arxiv.org/abs/2605.07587v1. Full licence notice: \texttt{licenses/arxiv-2605.07587-CC-BY-4.0.txt}.\par\smallskip

\noindent\textit{Editorial scope.} Counted unit: the article's theorem theo:recursiveAk (source0.tex lines 460--472). The article's complete original proof of it is delivered with it (source0.tex lines 474--519, one \texttt{proof} environment of 46 lines). No mathematics is written, repaired or re-derived: the statement and the proof are the article's own lines, in the article's own notation, and no retained line is substituted or re-flowed.  Retained uncounted as the source-local context the proof needs: lines 449--457.  One declared adaptation: exactly 1 ancillary strings inside the retained spans are omitted (image-inclusion commands and, where the source repeats a label, the repeated label definitions) and no other line differs from the source; the figure environments, with their placement, captions and labels, are kept, and the package records the omitted strings in fidelity receipts bound to the affected bodies.  No retained line contains a citation, so the wrapper carries no bibliography and no bibliography entry.  The retained lines contain the article's own diagram code, which the wrapper typesets with the standard tikz packages; no image file is delivered and the package declares no asset dependency.  Editorial formatting only, in addition: an AMS \texttt{amsart} preamble that restates the article's own declarations and macros used by the retained lines, namely the environments \texttt{theorem} on separate counters, together with the standard packages those lines need.  The retained environments print in this document as Theorem 1 in order of appearance, whereas the article prints the same items with the numbering of its own full document.  The complete native source of the article as distributed in the arXiv e-print for this exact version is bundled unchanged as \texttt{sources/200092/source0.tex}.
\begin{figure}
  \centering
  
\caption{Decomposition used to obtain $F_k(z,u)$ from $F_{k-1}(z,u)$. 
Cut after the $(k-1)$th $U_2$ step and at the first minimum between the $(k-1)$th and the $k$th $U_2$ step. 
The three parts correspond to $F_{k-1}(z,u)$, the operator $u^\ell \mapsto \frac{u^{\ell+1}-E(z)^{\ell+1}}{u-E(z)}$ (with $E(z)=zD(z)$), and a suffix $\Mup(z,u)$. 
The weight of the last $U_2$ step is implemented by applying $\frac{1}{z^{k-2}}\frac{d}{dz}z^{k-1}$ to the resulting generating function.}
  \label{fig:decomposition_Fk}
\end{figure}

\begin{theorem}\label{theo:recursiveAk}
    The generating function $F_{k}(z,u)$ of bicolored Dyck meanders with $k$ colored $U_2$ steps ending with a $U_2$ step where $z$ marks the number of steps and $u$ the final level satisfies for $k \geq 1$
    \begin{align*}
        F_{k}(z,u) 
            &= \frac{1}{z^{k-2}}\frac{d}{dz} \left( z^{k-1} \frac{u F_{k-1}(z,u) -  E(z) F_{k-1}(z,E(z))}{u -  E(z)} \Mup(z,u) \right),
    \end{align*}
    with $F_0(z) = 1$. 
    The generating function $B_k(z)$ is then for $k \geq 0$ given by
    \begin{align*}
        B_k(z) 
            &= F_{k}(z, E(z)) D(z).
    \end{align*}
\end{theorem}

\begin{proof}
For $k=1$, a bicolored meander counted by $F_1(z,u)$ is a Dyck meander followed by an up step,
where this last up step is colored ($U_2$) and carries weight equal to the path length.
 This is obtained from $\Mup(z,u)$ by the standard pointing operator
$z\frac{d}{dz}$, hence
\begin{align}\label{eq:F1}
    F_1(z,u) 
        = z \frac{d}{dz} \Mup(z,u)
        = \frac{z uE'(z)}{(1-uE(z))^2}.
\end{align}

Assume $k\ge 1$. We compute $F_{k}(z,u)$ from $F_{k-1}(z,u)$ as follows.
Given a path with $k$ steps $U_2$, we cut the path into three parts at the following two points (see Figure~\ref{fig:decomposition_Fk}):
after the $(k-1)$th step~$U_2$, and at the first minimum between the $(k-1)$th and $k$th step~$U_2$.
Now we discuss each part separately. 
The first part is enumerated by $F_{k-1}(z,u)$, a path from the origin to a level $\ell$ marked by the variable $u$.
The second part starts at level $\ell$ and then runs to a level $0 \leq i \leq \ell$ ending with a down step if it is nonempty. 
We will discuss it below.
The third part is then a meander that starts from level $i$ and ends with a step $U_2$.
Therefore, the third part is modeled by a factor $\Mup(z,u)$ in which the weight of the last step has to be adjusted.
This weight is the length of the path plus $k-1$ which can be computed by the operator $\frac{1}{z^{k-2}}\frac{d}{dz} z^{k-1} (\cdot)$ applied to the full generating function. 
Indeed, for any monomial~$z^n$, 
\begin{align*}
\frac{1}{z^{k-2}}\frac{d}{dz}\!\left(z^{k-1}z^n\right) = (n+k-1)\,z^n,
\end{align*}
so the operator multiplies each contribution by the path length plus $k-1$.

For the second part, we use a last-passage-decomposition by cutting each time the path attains a new lower level for the first time.
Between levels the path is a classical Dyck path associated with $D(z)$ and the first step to a lower level is a step $D$ associated with $z$. 
Thus, a path from level $\ell$ to level $i$ corresponds to $(zD(z))^{\ell-i}$.
Finally, we sum over all possible levels $0 \leq i \leq \ell$ of the minimum.
In terms of operators, we replace the weight $u^{\ell}$ of a path ending at level $\ell$ as follows (recall that $E(z) = z D(z)$):
\begin{align*}
    u^{\ell} 
        &\mapsto u^{\ell} + E(z)u^{\ell-1} + E(z)^2 u^{\ell-2} + \dots + E(z)^{\ell} 
        =  \frac{u^{\ell+1} -  E(z)^{\ell+1}}{u -  E(z)}.
\end{align*}
Finally, combining all three parts, this gives the claimed recurrence for $F_{k}(z,u)$.
In order to compute $B_{k}(z)$ we perform another first-passage decomposition at the end to let the walk run from altitude $\ell$ back to $0$. 
In terms of generating functions this gives
$
    B_k(z) = F_{k}(z, E(z)) D(z),
$
where the final factor $D(z)$ corresponds to a possibly empty Dyck path at the end of the walk.
Note that setting $F_{0}(z,u)=1$ this also gives the correct result for $k=0$ and $k=1$.
\end{proof}

\end{document}
